% ============================================================================
% OpenOrbitLink — RTL-SDR Ground Station Technical Reference
% Comprehensive Guide for Satellite Signal Reception
% Author: Ashutosh Kumar (OpenOrbitLink Project)
% Date: June 2026
% Location Context: Bengaluru, India
% ============================================================================

\documentclass[12pt,a4paper]{article}

\usepackage[utf8]{inputenc}
\usepackage[T1]{fontenc}
\usepackage{lmodern}
\usepackage[margin=2.5cm]{geometry}
\usepackage{graphicx}
\usepackage{booktabs}
\usepackage{longtable}
\usepackage{hyperref}
\usepackage{xcolor}
\usepackage{amsmath}
\usepackage{siunitx}
\usepackage{enumitem}
\usepackage{fancyhdr}
\usepackage{tikz}
\usepackage{pgfplots}
\pgfplotsset{compat=1.18}

\definecolor{loraBlue}{RGB}{0,180,216}
\definecolor{spaceBlack}{RGB}{10,14,33}
\definecolor{alertRed}{RGB}{233,69,96}

\hypersetup{
    colorlinks=true,
    linkcolor=loraBlue,
    urlcolor=loraBlue,
    citecolor=loraBlue
}

\pagestyle{fancy}
\fancyhf{}
\fancyhead[L]{\textbf{OpenOrbitLink — RTL-SDR Ground Station Guide}}
\fancyhead[R]{June 2026}
\fancyfoot[C]{\thepage}

\title{
    \vspace{-1cm}
    \textbf{\Huge RTL-SDR Ground Station} \\[0.5cm]
    \textbf{\Large Technical Reference for Satellite Signal Reception} \\[0.3cm]
    \large OpenOrbitLink Project — Comprehensive Buying \& Setup Guide \\[0.2cm]
    \normalsize Bengaluru, India | June 2026
}
\author{Ashutosh Kumar \\ \texttt{github.com/Ashutosh0x/OpenOrbitLink}}
\date{}

\begin{document}
\maketitle
\tableofcontents
\newpage

% ============================================================================
\section{Introduction: What is RTL-SDR?}
% ============================================================================

RTL-SDR (Realtek Software Defined Radio) is a low-cost USB dongle that turns any 
computer into a wideband radio receiver. Originally designed as a DVB-T television 
tuner, the RTL2832U chip inside was discovered to support raw I/Q sample output, 
enabling reception of virtually any radio signal between 500~kHz and 1.766~GHz.

\subsection{Why RTL-SDR for OpenOrbitLink?}

OpenOrbitLink operates on the \textbf{868~MHz ISM band} using LoRa modulation to 
communicate with Low Earth Orbit (LEO) satellites. The RTL-SDR V4 covers this 
frequency perfectly, enabling:

\begin{itemize}
    \item \textbf{Passive monitoring} of LoRa satellite downlinks
    \item \textbf{Spectrum analysis} for link budget validation
    \item \textbf{Weather satellite} image reception (NOAA, Meteor-M2)
    \item \textbf{ADS-B} aircraft tracking (1090~MHz)
    \item \textbf{APRS} packet reception for mesh fallback
    \item \textbf{Signal quality} measurement and ground-truth validation
\end{itemize}

% ============================================================================
\section{Radio Wave Fundamentals}
% ============================================================================

\subsection{Electromagnetic Spectrum — Bands We Use}

\begin{table}[h!]
\centering
\caption{Radio Bands Relevant to OpenOrbitLink}
\begin{tabular}{@{}llll@{}}
\toprule
\textbf{Band} & \textbf{Frequency} & \textbf{Wavelength} & \textbf{Use Case} \\
\midrule
VHF & 30--300~MHz & 10--1~m & NOAA weather sats (137~MHz) \\
    &             &         & APRS (144.39~MHz) \\
    &             &         & ISS voice (145.8~MHz) \\
UHF & 300~MHz--3~GHz & 1~m--10~cm & \textbf{LoRa ISM (868~MHz)} \\
    &                &            & ADS-B aircraft (1090~MHz) \\
    &                &            & GPS L1 (1575.42~MHz) \\
L-Band & 1--2~GHz & 30--15~cm & Iridium (1616--1626~MHz) \\
S-Band & 2--4~GHz & 15--7.5~cm & Amateur satellites \\
\bottomrule
\end{tabular}
\end{table}

\subsection{Free-Space Path Loss (FSPL)}

The signal weakens as it travels from satellite to ground. The loss is calculated as:

\begin{equation}
    \text{FSPL (dB)} = 20 \log_{10}(d) + 20 \log_{10}(f) + 20 \log_{10}\left(\frac{4\pi}{c}\right)
\end{equation}

Where $d$ = distance in meters, $f$ = frequency in Hz, $c$ = speed of light.

\begin{table}[h!]
\centering
\caption{FSPL at 868~MHz for Different Distances}
\begin{tabular}{@{}lll@{}}
\toprule
\textbf{Scenario} & \textbf{Distance} & \textbf{FSPL (dB)} \\
\midrule
Ground-to-Ground (rural) & 15~km & 111.3 \\
Ground-to-Ground (record) & 832~km & 147.6 \\
Satellite overhead (zenith) & 550~km & 144.1 \\
Satellite at horizon (5°) & 2,100~km & 155.7 \\
Satellite at 20° elevation & 1,200~km & 150.8 \\
Satellite at 45° elevation & 700~km & 146.1 \\
\bottomrule
\end{tabular}
\end{table}

\subsection{LoRa Modulation — Chirp Spread Spectrum}

LoRa uses \textbf{Chirp Spread Spectrum (CSS)} modulation, where each symbol is a 
frequency sweep (chirp) across the bandwidth. This provides exceptional sensitivity:

\begin{table}[h!]
\centering
\caption{LoRa Spreading Factors and Performance at BW=125~kHz}
\begin{tabular}{@{}llllll@{}}
\toprule
\textbf{SF} & \textbf{Bitrate} & \textbf{Sensitivity} & \textbf{SNR Limit} & \textbf{Range (Sat)} & \textbf{ToA (51B)} \\
\midrule
SF7  & 5,469~bps & $-123$~dBm & $-7.5$~dB & Overhead only & 92~ms \\
SF8  & 3,125~bps & $-126$~dBm & $-10$~dB  & $>45°$ elev & 165~ms \\
SF9  & 1,758~bps & $-129$~dBm & $-12.5$~dB & $>30°$ elev & 329~ms \\
SF10 & 977~bps   & $-132$~dBm & $-15$~dB  & $>20°$ elev & 657~ms \\
SF11 & 537~bps   & $-135$~dBm & $-17.5$~dB & $>10°$ elev & 1,315~ms \\
SF12 & 293~bps   & $-137$~dBm & $-20$~dB  & Horizon & 2,466~ms \\
\bottomrule
\end{tabular}
\end{table}

% ============================================================================
\section{Hardware: RTL-SDR Blog V4 — Complete Specifications}
% ============================================================================

\subsection{Core Specifications}

\begin{table}[h!]
\centering
\caption{RTL-SDR Blog V4 Technical Specifications}
\begin{tabular}{@{}ll@{}}
\toprule
\textbf{Parameter} & \textbf{Value} \\
\midrule
ADC Chip & Realtek RTL2832U \\
Tuner Chip & Rafael Micro R828D \\
Frequency Range & 500~kHz -- 1,766~MHz \\
Maximum Bandwidth & 2.4~MHz (stable), 3.2~MHz (max) \\
ADC Resolution & 8-bit I/Q \\
Sample Rate & Up to 2.4~MSPS (stable) \\
TCXO Accuracy & $\pm 1$~PPM (temperature-compensated) \\
Input Connector & SMA Female (50$\Omega$) \\
Bias-Tee & Software-switchable, 4.5~V, 180~mA max \\
HF Mode & Built-in upconverter (no aliasing) \\
Input Filtering & Triplexed: HF / VHF / UHF paths \\
Notch Filters & FM broadcast \& DAB rejection \\
Case & Aluminum shielded enclosure \\
Power & USB powered ($\sim$280~mA at 5~V) \\
Interface & USB 2.0 \\
Dimensions & 65 $\times$ 30 $\times$ 15~mm \\
Weight & $\sim$50~g \\
Operating Temp & $-10°$C to $+60°$C \\
\bottomrule
\end{tabular}
\end{table}

\subsection{Sensitivity and Dynamic Range}

\begin{table}[h!]
\centering
\caption{Measured Sensitivity at Key Frequencies}
\begin{tabular}{@{}lll@{}}
\toprule
\textbf{Frequency} & \textbf{Noise Floor} & \textbf{Application} \\
\midrule
137~MHz & $-142$~dBm & NOAA weather satellite APT \\
144~MHz & $-141$~dBm & APRS, amateur radio \\
435~MHz & $-140$~dBm & Amateur satellite downlinks \\
\textbf{868~MHz} & \textbf{$-138$~dBm} & \textbf{LoRa ISM (OpenOrbitLink)} \\
1090~MHz & $-136$~dBm & ADS-B aircraft tracking \\
1575~MHz & $-133$~dBm & GPS L1 (near upper limit) \\
\bottomrule
\end{tabular}
\end{table}

% ============================================================================
\section{Reception Range and Distance Analysis}
% ============================================================================

\subsection{Maximum Reception Distances by Signal Type}

\begin{table}[h!]
\centering
\caption{Achievable Reception Distances with RTL-SDR V4}
\begin{tabular}{@{}lllp{4.5cm}@{}}
\toprule
\textbf{Signal} & \textbf{Freq} & \textbf{Max Range} & \textbf{Notes} \\
\midrule
FM Radio & 88--108~MHz & 100--300~km & Line of sight + diffraction \\
NOAA APT & 137~MHz & 2,800~km & Satellite footprint horizon-to-horizon \\
Meteor LRPT & 137~MHz & 2,800~km & Higher quality images \\
APRS & 144~MHz & 50--100~km & Ground-to-ground, terrain dependent \\
ISS Voice & 145.8~MHz & 2,200~km & During overhead pass \\
ADS-B & 1090~MHz & 300--450~km & Aircraft at 10,000~m altitude \\
\textbf{LoRa Sat} & \textbf{868~MHz} & \textbf{2,100~km} & \textbf{Satellite at 5° elevation (SF12)} \\
LoRa Ground & 868~MHz & 15--50~km & Rural, clear line of sight \\
LoRa Record & 868~MHz & 832~km & Ground-to-ground world record \\
\bottomrule
\end{tabular}
\end{table}

\subsection{Link Budget Calculation for LoRa Satellite}

Complete link budget for FOSSASAT-2E to RTL-SDR ground station in Bengaluru:

\begin{align}
    P_{\text{TX}} &= +20 \text{~dBm} && \text{(Satellite TX power: 100~mW)} \\
    G_{\text{TX}} &= +2 \text{~dBi} && \text{(Satellite antenna: omnidirectional)} \\
    L_{\text{FSPL}} &= -154 \text{~dB} && \text{(868~MHz, 700~km slant range)} \\
    L_{\text{atm}} &= -1.5 \text{~dB} && \text{(Atmospheric absorption)} \\
    L_{\text{pol}} &= -3 \text{~dB} && \text{(Polarization mismatch)} \\
    L_{\text{misc}} &= -1 \text{~dB} && \text{(Cable, connector losses)} \\
    G_{\text{RX}} &= +5 \text{~dBi} && \text{(Ground Yagi antenna)} \\
    G_{\text{LNA}} &= +20 \text{~dB} && \text{(Low Noise Amplifier)} \\
    \hline
    P_{\text{RX}} &= -112.5 \text{~dBm} && \text{(Signal at SDR input)} \\
    N_{\text{floor}} &= -138 \text{~dBm} && \text{(RTL-SDR noise floor)} \\
    \text{SNR} &= +25.5 \text{~dB} && \text{(Excellent — link closes easily)}
\end{align}

\textbf{Without LNA} (RTL-SDR + Yagi only):
\begin{align}
    P_{\text{RX}} &= -132.5 \text{~dBm} \\
    \text{SNR} &= +5.5 \text{~dB} && \text{(Marginal but works with SF12)}
\end{align}

\textbf{With dipole antenna only} (RTL-SDR kit antenna):
\begin{align}
    P_{\text{RX}} &= -137.5 \text{~dBm} \\
    \text{SNR} &= +0.5 \text{~dB} && \text{(Barely possible — LNA strongly recommended)}
\end{align}

% ============================================================================
\section{Complete Hardware Shopping List}
% ============================================================================

\subsection{Essential Components}

\begin{longtable}{@{}p{4cm}p{2cm}p{2cm}p{4.5cm}@{}}
\caption{Complete Shopping List — Bengaluru, India (June 2026)} \\
\toprule
\textbf{Component} & \textbf{Price (\rupee)} & \textbf{Priority} & \textbf{Where to Buy} \\
\midrule
\endfirsthead
\toprule
\textbf{Component} & \textbf{Price (\rupee)} & \textbf{Priority} & \textbf{Where to Buy} \\
\midrule
\endhead

\multicolumn{4}{l}{\textbf{Core SDR Hardware}} \\
\midrule
RTL-SDR Blog V4 Kit (with dipole antennas) & 6,000--7,500 & Essential & Robu.in, FabToLab.com, ElectroPi.in \\
USB Extension Cable (3m, shielded) & 200--400 & Essential & Amazon.in, local shops \\
SMA Male-Male adapter cable & 150--300 & Essential & Robu.in, SP Road (Blr) \\

\midrule
\multicolumn{4}{l}{\textbf{Antennas}} \\
\midrule
868~MHz Quarter-wave ground plane (DIY) & 100--300 & Essential & DIY with copper wire \\
868~MHz 5-element Yagi (12~dBi) & 900--2,000 & Recommended & Eteily Tech, IndiaMART \\
137~MHz V-Dipole (for weather sats) & 0 (DIY) & Optional & Built from kit dipoles \\
QFH Antenna (omnidirectional, for sats) & 500--1,500 & Optional & DIY or IndiaMART \\

\midrule
\multicolumn{4}{l}{\textbf{Amplification}} \\
\midrule
LNA (SPF5189Z wideband preamp) & 800--1,200 & Recommended & Robu.in, AliExpress \\
RTL-SDR Blog Wideband LNA & 2,000--2,500 & Best option & FabToLab.com \\
Bandpass Filter 868~MHz (SAW) & 400--800 & Optional & ElectronicsComp \\

\midrule
\multicolumn{4}{l}{\textbf{LoRa Transmission Module}} \\
\midrule
SX1276 LoRa Module (868~MHz) & 450--800 & Essential & Robu.in, SP Road \\
SX1262 LoRa Module (advanced) & 600--1,200 & Upgrade & ElectronicsComp \\
E32-868T30D (1W LoRa module) & 1,000--1,500 & High power & IndiaMART, Quartz \\
LoRa antenna (SMA, 868~MHz) & 200--400 & Essential & Robu.in \\

\midrule
\multicolumn{4}{l}{\textbf{Computing}} \\
\midrule
Raspberry Pi 4 (4GB) & 4,000--5,000 & For 24/7 station & Robu.in, CanaKit.in \\
MicroSD Card (32GB) & 400--600 & With RPi & Amazon.in \\
Raspberry Pi case + fan & 500--800 & With RPi & Robu.in \\

\midrule
\multicolumn{4}{l}{\textbf{Cables \& Connectors}} \\
\midrule
RG58 Coax cable (5m) & 300--500 & Essential & SP Road, IndiaMART \\
SMA to BNC adapter & 100--200 & Useful & Electronics shops \\
SMA to N-Type adapter & 150--250 & For Yagi & IndiaMART \\
Weatherproof enclosure (IP65) & 400--800 & Outdoor setup & Amazon.in \\

\midrule
\textbf{TOTAL (Minimum Setup)} & \textbf{$\sim$7,800} & & \\
\textbf{TOTAL (Full Station)} & \textbf{$\sim$22,000} & & \\
\bottomrule
\end{longtable}

\subsection{Where to Buy in Bengaluru}

\begin{table}[h!]
\centering
\caption{Bengaluru Electronics Shopping Guide}
\begin{tabular}{@{}llp{5cm}@{}}
\toprule
\textbf{Store / Area} & \textbf{Type} & \textbf{What to Find} \\
\midrule
\textbf{SP Road} & Physical & Connectors, cables, SMA adapters, wire for DIY antennas \\
\textbf{National Market (Chickpete)} & Physical & Coax cable, enclosures, general electronics \\
\textbf{Robu.in} & Online & RTL-SDR, LoRa modules, Raspberry Pi \\
\textbf{FabToLab.com} & Online & RTL-SDR Blog products, LNAs \\
\textbf{ElectroPi.in} & Online & SDR accessories, dev boards \\
\textbf{ElectronicsComp.com} & Online & SX1276/SX1262 modules, filters \\
\textbf{Sharvielectronics.com} & Online & LoRa modules, antennas \\
\textbf{GSAS Micro Systems} & Bangalore & Authorized Reyax LoRa partner \\
\textbf{IndiaMART} & B2B & Yagi antennas, bulk components \\
\bottomrule
\end{tabular}
\end{table}

% ============================================================================
\section{Signal Reception Capabilities}
% ============================================================================

\subsection{What You CAN Receive}

\begin{longtable}{@{}p{3.5cm}lllp{3.5cm}@{}}
\caption{Complete List of Receivable Signals} \\
\toprule
\textbf{Signal} & \textbf{Freq (MHz)} & \textbf{Range} & \textbf{Antenna} & \textbf{Software} \\
\midrule
\endfirsthead

FM Broadcast & 88--108 & 100~km & Kit dipole & SDR\#, GQRX \\
NOAA-18 APT & 137.9125 & 2,800~km & V-dipole & WXtoImg, noaa-apt \\
NOAA-19 APT & 137.1000 & 2,800~km & V-dipole & WXtoImg, noaa-apt \\
Meteor-M2 LRPT & 137.1000 & 2,800~km & V-dipole/QFH & MeteorDemod \\
APRS & 144.3900 & 80~km & Dipole & Direwolf \\
ISS APRS & 145.8250 & 2,200~km & Dipole & Direwolf \\
ISS Voice & 145.8000 & 2,200~km & Dipole & SDR\# (FM) \\
\textbf{LoRa IoT Sats} & \textbf{868.1} & \textbf{2,100~km} & \textbf{Yagi} & \textbf{GNU Radio, gr-lora} \\
\textbf{FOSSASAT} & \textbf{868.0} & \textbf{2,100~km} & \textbf{Yagi} & \textbf{FOSSA Ground Station} \\
\textbf{Lacuna Space} & \textbf{868.1} & \textbf{2,100~km} & \textbf{Yagi} & \textbf{TinyGS} \\
ADS-B (aircraft) & 1090 & 450~km & 1090 ant & dump1090, tar1090 \\
Iridium & 1616--1626 & 2,000~km & Patch & iridium-toolkit \\
GPS L1 & 1575.42 & 20,200~km & Patch & gnss-sdr \\
\bottomrule
\end{longtable}

\subsection{What You CANNOT Receive}

\begin{table}[h!]
\centering
\caption{Signals Beyond RTL-SDR Capability}
\begin{tabular}{@{}llp{5cm}@{}}
\toprule
\textbf{Signal} & \textbf{Frequency} & \textbf{Why Not} \\
\midrule
Starlink internet & 10.7--12.7~GHz & Ku-band — far above RTL-SDR range \\
Starlink uplink & 14.0--14.5~GHz & Ka-band — needs dish + LNB \\
5G mmWave & 24--39~GHz & Millimeter wave — specialized hardware \\
Military SATCOM & Various & Encrypted + classified frequencies \\
Cell phone calls & 700--2600~MHz & Encrypted (illegal to intercept) \\
WiFi & 2.4/5~GHz & Above RTL-SDR range (partially) \\
\bottomrule
\end{tabular}
\end{table}

% ============================================================================
\section{The Complete Signal Chain}
% ============================================================================

\subsection{End-to-End Signal Flow}

The following describes how a LoRa satellite signal travels from orbit to your 
OpenOrbitLink app:

\begin{enumerate}
    \item \textbf{Satellite Transmitter} — FOSSASAT-2E at 550~km altitude broadcasts
          a LoRa packet at 868~MHz with 100~mW (+20~dBm) through an omnidirectional
          antenna (+2~dBi gain).
    
    \item \textbf{Propagation} — Signal travels 550--2,100~km through vacuum and 
          atmosphere, losing 144--156~dB of power (Free-Space Path Loss + atmospheric 
          absorption + polarization mismatch).
    
    \item \textbf{Ground Antenna} — Your 868~MHz Yagi antenna captures the 
          electromagnetic wave and converts it to an electrical signal of approximately 
          $-130$~dBm ($10^{-16}$~watts).
    
    \item \textbf{LNA (Optional)} — Low Noise Amplifier boosts the signal by +20~dB 
          while adding minimal noise (NF $< 1$~dB), bringing signal to $-110$~dBm.
    
    \item \textbf{RTL-SDR V4 — R828D Tuner} — The tuner chip selects exactly 868.1~MHz,
          rejects out-of-band interference, and downconverts the RF signal to an 
          intermediate frequency (IF).
    
    \item \textbf{RTL-SDR V4 — RTL2832U ADC} — The analog-to-digital converter samples 
          the IF signal at 2.4~million samples/second with 8-bit resolution, producing 
          raw I/Q (In-phase/Quadrature) data.
    
    \item \textbf{USB Transfer} — Raw I/Q samples stream over USB 2.0 to your computer 
          at approximately 4.8~MB/s.
    
    \item \textbf{SDR Software} — SDR\# or GQRX displays the waterfall spectrogram.
          LoRa chirps appear as diagonal lines sweeping across the bandwidth.
    
    \item \textbf{LoRa Demodulator} — GNU Radio with \texttt{gr-lora} block, or 
          dedicated LoRa decoder extracts the raw packet data from CSS chirps.
    
    \item \textbf{OpenOrbitLink Backend} — Python backend processes the BPv7 DTN 
          bundle, decrypts payload, validates integrity.
    
    \item \textbf{Android App} — Your phone displays the received message, satellite 
          telemetry, or queued response.
\end{enumerate}

% ============================================================================
\section{Software Stack}
% ============================================================================

\begin{table}[h!]
\centering
\caption{Complete Software Stack for Ground Station}
\begin{tabular}{@{}lllp{3.5cm}@{}}
\toprule
\textbf{Software} & \textbf{Platform} & \textbf{License} & \textbf{Purpose} \\
\midrule
\multicolumn{4}{l}{\textbf{SDR Interface}} \\
\midrule
SDR\# (SDRSharp) & Windows & Freeware & General SDR receiver \\
GQRX & Linux/Mac & Open source & SDR receiver + waterfall \\
CubicSDR & All & Open source & Cross-platform SDR \\
SDRangel & All & Open source & Advanced multi-channel \\

\midrule
\multicolumn{4}{l}{\textbf{Satellite Tracking}} \\
\midrule
Gpredict & All & Open source & Pass prediction + Doppler \\
SatDump & All & Open source & Satellite signal decoder \\
Orbitron & Windows & Freeware & TLE-based tracking \\

\midrule
\multicolumn{4}{l}{\textbf{LoRa Decoding}} \\
\midrule
GNU Radio & Linux/Win & Open source & Signal processing framework \\
gr-lora & Linux & Open source & LoRa demodulator block \\
TinyGS & ESP32 & Open source & LoRa ground station \\
FOSSA GS & All & Open source & FOSSASAT decoder \\

\midrule
\multicolumn{4}{l}{\textbf{Weather Satellites}} \\
\midrule
WXtoImg & Windows & Freeware & NOAA APT image decoder \\
MeteorDemod & All & Open source & Meteor-M2 LRPT decoder \\
noaa-apt & All & Open source & Rust-based APT decoder \\
SatDump & All & Open source & Universal satellite decoder \\

\midrule
\multicolumn{4}{l}{\textbf{Aircraft Tracking}} \\
\midrule
dump1090 & All & Open source & ADS-B decoder \\
tar1090 & Linux & Open source & Web-based ADS-B display \\
FlightAware & Cloud & Free tier & ADS-B feed network \\

\midrule
\multicolumn{4}{l}{\textbf{Drivers}} \\
\midrule
Zadig & Windows & Open source & USB driver installer \\
rtl-sdr (librtlsdr) & Linux & Open source & RTL-SDR driver library \\
\bottomrule
\end{tabular}
\end{table}

% ============================================================================
\section{Comparison: RTL-SDR vs Alternatives}
% ============================================================================

\begin{table}[h!]
\centering
\caption{SDR Hardware Comparison (June 2026 Pricing)}
\begin{tabular}{@{}lllll@{}}
\toprule
\textbf{Feature} & \textbf{RTL-SDR V4} & \textbf{Nooelec V5} & \textbf{HackRF One} & \textbf{Airspy Mini} \\
\midrule
Price (India) & \rupee 6,000 & \rupee 5,000 & \rupee 35,000 & \rupee 15,000 \\
TX Capable & No & No & \textbf{Yes} & No \\
Freq Range & 500k--1.7G & 25M--1.7G & 1M--6G & 24M--1.8G \\
Bandwidth & 2.4~MHz & 2.4~MHz & 20~MHz & 6~MHz \\
ADC Bits & 8-bit & 8-bit & 8-bit & 12-bit \\
TCXO & 1~PPM & 0.5~PPM & 20~PPM & 0.5~PPM \\
Bias-Tee & Software & Always-on & Software & Software \\
HF Support & Upconverter & Direct samp & Native & No \\
Best For & \textbf{Value + LoRa} & Monitoring & TX + wide BW & Sensitivity \\
\bottomrule
\end{tabular}
\end{table}

\subsection{Verdict for OpenOrbitLink}

\textbf{RTL-SDR Blog V4} is the recommended choice because:
\begin{enumerate}
    \item Covers 868~MHz LoRa band with excellent sensitivity ($-138$~dBm noise floor)
    \item Software-switchable Bias-Tee safely powers LNA without hardware modification
    \item 1~PPM TCXO provides frequency stability critical for LoRa demodulation
    \item Best price-to-performance ratio at \rupee 6,000
    \item No TX needed — LoRa transmission is handled by the SX1276/SX1262 module
\end{enumerate}

% ============================================================================
\section{Cost Comparison: OpenOrbitLink vs Starlink}
% ============================================================================

\begin{table}[h!]
\centering
\caption{Total Cost of Ownership — OpenOrbitLink vs Starlink}
\begin{tabular}{@{}lll@{}}
\toprule
\textbf{Item} & \textbf{OpenOrbitLink} & \textbf{Starlink} \\
\midrule
Hardware (one-time) & \rupee 7,800 (\$93) & \$599 (\rupee 50,000) \\
Monthly subscription & \rupee 0 (free) & \$120/month (\rupee 10,000) \\
Year 1 total & \rupee 7,800 & \rupee 1,70,000 \\
Year 3 total & \rupee 7,800 & \rupee 4,10,000 \\
License required & No (ISM band) & No \\
Internet access & No (messaging only) & Yes (broadband) \\
Power consumption & 2W (USB) & 75W (dish + router) \\
Works off-grid & Yes (battery) & No (needs 75W) \\
\bottomrule
\end{tabular}
\end{table}

\vspace{0.5cm}
\textbf{OpenOrbitLink is 52$\times$ cheaper over 3 years.} While it provides messaging 
rather than broadband internet, it serves the critical need of \textbf{connectivity in 
zero-infrastructure areas} where Starlink is impractical due to power and cost.

% ============================================================================
\section{Quick Start: Your First Satellite Reception}
% ============================================================================

\subsection{Step 1: Unbox and Install (15 minutes)}
\begin{enumerate}
    \item Connect RTL-SDR V4 to USB extension cable
    \item Attach the kit's telescopic dipole antenna
    \item Install Zadig driver (Windows) or \texttt{apt install rtl-sdr} (Linux)
    \item Install SDR\# or GQRX
    \item Tune to local FM station (e.g., 93.5~MHz) to verify operation
\end{enumerate}

\subsection{Step 2: Receive NOAA Weather Satellite (Day 1)}
\begin{enumerate}
    \item Open Gpredict, add NOAA-19 (NORAD ID: 33591)
    \item Set your location to Bengaluru (12.97°N, 77.59°E)
    \item Wait for next pass ($>30°$ elevation recommended)
    \item Set SDR\# to 137.100~MHz, FM Wide, 40~kHz bandwidth
    \item Arrange dipole as V-shape, pointing North-South
    \item Record audio during pass ($\sim$12 minutes)
    \item Decode with WXtoImg $\rightarrow$ weather image from space!
\end{enumerate}

\subsection{Step 3: Monitor LoRa Satellites (Week 1)}
\begin{enumerate}
    \item Build or buy 868~MHz Yagi antenna
    \item Install GNU Radio + \texttt{gr-lora} blocks
    \item Track FOSSASAT-2E (NORAD ID: TBD) via CelesTrak
    \item Tune to 868.0~MHz during predicted pass
    \item Look for chirp patterns on waterfall display
    \item Decode LoRa packets with gr-lora
\end{enumerate}

% ============================================================================
\section{Useful Links}
% ============================================================================

\begin{itemize}
    \item \textbf{RTL-SDR Blog:} \url{https://www.rtl-sdr.com}
    \item \textbf{Robu.in (Buy):} \url{https://robu.in}
    \item \textbf{FabToLab (Buy):} \url{https://fabtolab.com}
    \item \textbf{ElectroPi (Buy):} \url{https://electropi.in}
    \item \textbf{CompoIndia (Buy):} \url{https://compoindia.com}
    \item \textbf{The Engineer Store (Bengaluru):} \url{https://theengineerstore.in}
    \item \textbf{VirtusFab (Buy):} \url{https://virtusfab.in}
    \item \textbf{MG Super Labs (Buy):} \url{https://mgsuperlabs.co.in}
    \item \textbf{element14 India:} \url{https://element14.com/india}
    \item \textbf{Mouser India:} \url{https://mouser.in}
    \item \textbf{CelesTrak (TLE data):} \url{https://celestrak.org}
    \item \textbf{TinyGS (LoRa ground station):} \url{https://tinygs.com}
    \item \textbf{FOSSA Systems:} \url{https://fossa.systems}
    \item \textbf{GNU Radio:} \url{https://www.gnuradio.org}
    \item \textbf{gr-lora (LoRa decoder):} \url{https://github.com/rpp0/gr-lora}
    \item \textbf{SatDump:} \url{https://github.com/SatDump/SatDump}
    \item \textbf{Gpredict:} \url{http://gpredict.oz9aec.net}
    \item \textbf{r/RTLSDR:} \url{https://reddit.com/r/RTLSDR}
    \item \textbf{OpenOrbitLink:} \url{https://github.com/Ashutosh0x/OpenOrbitLink}
\end{itemize}

% ============================================================================
\section{Advanced SDR Hardware — Beyond RTL-SDR}
% ============================================================================

For users who need \textbf{transmit capability}, higher bandwidth, or full-duplex 
operation, the following SDR platforms offer significantly more power than the RTL-SDR V4.

\subsection{HackRF One V1.9.1 — Best Value TX/RX SDR}

The \textbf{HackRF One} is the most popular transmit-capable SDR, covering an 
extraordinary 1~MHz to 6~GHz range. The V1.9.1 (Clifford Heath version) includes 
major hardware improvements over earlier revisions.

\begin{table}[h!]
\centering
\caption{HackRF One V1.9.1 Complete Specifications}
\begin{tabular}{@{}ll@{}}
\toprule
\textbf{Parameter} & \textbf{Value} \\
\midrule
Frequency Range & 1~MHz -- 6~GHz \\
Bandwidth & Up to 20~MHz \\
Sample Rate & Up to 20~MSPS \\
ADC Resolution & 8-bit I/Q \\
Duplex & Half-duplex (TX or RX, not simultaneous) \\
TX Power & Up to +15~dBm (30~mW) \\
RX LNA & TRF37B73 MMIC (V1.9.1 upgrade) \\
Antenna Connector & SMA female (50$\Omega$) \\
Bias-Tee & Software-controlled, 50~mA at 3.3~V \\
Clock & Internal; external TCXO optional \\
Clock I/O & SMA connectors for sync \\
Antenna Protection & SMP1330-085LF (V1.9.1 upgrade) \\
RF Switches & SKY13453-385LF (current production) \\
USB & Hi-Speed USB 2.0 \\
Power & USB-powered \\
Case & Aluminum (optional) \\
Open Source & Yes --- hardware and firmware \\
\bottomrule
\end{tabular}
\end{table}

\subsubsection{V1.9.1 Hardware Improvements}

The Clifford Heath version (V1.9.1) includes critical upgrades:

\begin{enumerate}
    \item \textbf{Antenna port protection:} SMP1330-085LF protects both TX amplifier 
          and RX LNA from excessive RF input
    \item \textbf{Improved Bias-T:} Operates at higher frequencies, improving RF 
          sensitivity whether enabled or not
    \item \textbf{Updated RF switches:} SKY13453-385LF (current production, replacing 
          discontinued SKY13350-385LF)
    \item \textbf{New LNA:} TRF37B73 replaces MGA-81563-TR1G
    \item \textbf{Updated USB protection:} USB6B1 (current production)
    \item \textbf{Updated clock protection:} BAS70 diodes on CLKIN/CLKOUT
    \item \textbf{Reduced PCB drill sizes:} 14 down to 6, all vias 27/13 mil
\end{enumerate}

\subsubsection{FabToLab India Pricing (June 2026, Verified In-Stock)}

\textbf{Store:} Fab.to.Lab (\url{https://fabtolab.com}) \\
\textbf{Phone:} +91-9110845041, +91-8050032228/9 \\
\textbf{Product Code:} PRC-HackRF-One \\
\textbf{Status:} \textcolor{green!60!black}{\textbf{IN STOCK}} \\

\begin{table}[h!]
\centering
\caption{HackRF One Pricing on FabToLab (June 2026)}
\begin{tabular}{@{}clll@{}}
\toprule
\textbf{Option} & \textbf{Description} & \textbf{Price} & \textbf{Includes} \\
\midrule
82 & Single motherboard & \rupee 19,502 & HackRF board only \\
83 & Board + case & \rupee 20,431 & + Aluminum enclosure \\
84 & Full Set & \rupee 24,183 & + Case + antennas + USB \\
\textbf{85} & \textbf{Full + TCXO 0.1PPM} & \textbf{\rupee 25,744} & \textbf{+ Precision oscillator} \\
\bottomrule
\end{tabular}
\end{table}

\textbf{Bulk discounts:} 2+ units: \rupee 19,346 each | 4+ units: \rupee 19,188 each | 7+ units: \rupee 19,031 each

\vspace{0.3cm}
\noindent\fbox{\parbox{\textwidth}{
\textbf{Recommendation:} Choose \textbf{Option 85} (\rupee 25,744). The TCXO 0.1~PPM 
upgrade is critical --- standard HackRF has 20~PPM drift, meaning at 868~MHz the 
frequency wanders by $\pm$17~kHz. With 0.1~PPM, drift is only $\pm$87~Hz. LoRa 
demodulation requires this precision for reliable packet decoding.
}}

\subsubsection{HackRF Accessories on FabToLab}

\begin{table}[h!]
\centering
\caption{HackRF Accessories Available on FabToLab (In-Stock)}
\begin{tabular}{@{}llll@{}}
\toprule
\textbf{Accessory} & \textbf{Price} & \textbf{Covers} & \textbf{Need?} \\
\midrule
PortaPack H4M & \rupee 29,503 & Standalone touchscreen & Highly recommended \\
HackRF Pro (newest!) & \rupee 30,495 & Next-gen HackRF & Alternative to V1.9.1 \\
ANT500 telescope antenna & \rupee 2,846 & 75~MHz -- 1~GHz & Good general use \\
ANT700 telescope antenna & \rupee 2,372 & 300~MHz -- 1.1~GHz & Covers 868~MHz! \\
800--2200~MHz antenna & \rupee 3,140 & LoRa + ADS-B & Best for satellites \\
Aluminum Enclosure Kit & \rupee 3,353 & RF shielding & If buying bare board \\
RF Shield + Component Kit & \rupee 1,336 & Internal shielding & Optional \\
Opera Cake (antenna switch) & \rupee 16,941 & Multi-antenna auto-switch & Advanced \\
Tiny TCXO 0.5~PPM module & \rupee 2,949 & Frequency stability & If not option 85 \\
ESP32-MDK (WiFi addon) & \rupee 7,420 & WiFi for PortaPack & Optional IoT \\
Ham It Up v1.3 (HF upconverter) & \rupee 6,040 & Below 30~MHz & Shortwave only \\
Acrylic Case & \rupee 1,423 & Transparent enclosure & Budget option \\
\bottomrule
\end{tabular}
\end{table}

\subsection{Other SDR Platforms Available in India}

\begin{table}[h!]
\centering
\caption{Advanced SDR Comparison with India Availability (June 2026)}
\resizebox{\textwidth}{!}{
\begin{tabular}{@{}lllllllll@{}}
\toprule
\textbf{SDR} & \textbf{Freq} & \textbf{BW} & \textbf{TX} & \textbf{Duplex} & \textbf{ADC} & \textbf{Price (\rupee)} & \textbf{Stock} & \textbf{Buy From} \\
\midrule
RTL-SDR V4 & 500k--1.7G & 2.4M & No & --- & 8-bit & 6,000 & Scarce & FabToLab \\
SDRplay RSP1B & 1k--2G & 10M & No & --- & 14-bit & 12,000 & Yes & FabToLab \\
Airspy Mini & 24M--1.8G & 6M & No & --- & 12-bit & 15,000 & Limited & Mouser.in \\
\textbf{HackRF One} & \textbf{1M--6G} & \textbf{20M} & \textbf{Yes} & Half & 8-bit & \textbf{19,502} & \textbf{Yes} & \textbf{FabToLab} \\
HackRF Pro & 1M--6G & 20M & Yes & Half & 8-bit & 30,495 & Yes & FabToLab \\
ADALM-Pluto & 325M--3.8G & 20M & Yes & Full & 12-bit & 25,000 & Yes & element14 \\
LimeSDR Mini 2 & 10M--3.5G & 30.7M & Yes & Full & 12-bit & 40,000 & Yes & EngineerStore \\
HydraSDR RFOne & 24M--1.8G & 10M & No & --- & 12-bit & 27,093 & Yes & FabToLab \\
Fobos SDR & 50M--6G & 50M & No & USB3 & 14-bit & 53,680 & Yes & FabToLab \\
SignalSDR Pro & 1M--6G & 61M & Yes & Full & 12-bit & 1,61,317 & Yes & FabToLab \\
BladeRF 2.0 xA9 & 47M--6G & 61.4M & Yes & Full & 12-bit & 75,000 & Import & Mouser.in \\
USRP B210 & 70M--6G & 56M & Yes & Full & 12-bit & 2,50,000 & Import & Mouser.in \\
\bottomrule
\end{tabular}
}
\end{table}

\subsection{Where to Buy SDR Hardware in India}

\begin{table}[h!]
\centering
\caption{Indian SDR Retailers — Verified June 2026}
\begin{tabular}{@{}lllp{4cm}@{}}
\toprule
\textbf{Store} & \textbf{Location} & \textbf{URL} & \textbf{Best For} \\
\midrule
\textbf{FabToLab} & Bengaluru & \url{fabtolab.com} & HackRF, accessories, widest SDR selection \\
\textbf{The Engineer Store} & \textbf{Bengaluru} & \url{theengineerstore.in} & LimeSDR Mini 2.0 (local warehouse!) \\
\textbf{CompoIndia} & Delhi & \url{compoindia.com} & HackRF + PortaPack kits \\
\textbf{ElectroPi} & Delhi & \url{electropi.in} & Budget SDR, LoRa modules \\
\textbf{Robu.in} & Pune & \url{robu.in} & RTL-SDR, LoRa, Raspberry Pi \\
\textbf{VirtusFab} & Online & \url{virtusfab.in} & HackRF + PortaPack bundles \\
\textbf{MG Super Labs} & Mumbai & \url{mgsuperlabs.co.in} & USRP, BladeRF, pro SDRs \\
\textbf{element14 India} & Online & \url{element14.com/india} & ADALM-Pluto (authorized Analog Devices) \\
\textbf{Mouser India} & Online & \url{mouser.in} & All brands, no customs hassle \\
\textbf{IndiaMART} & B2B & \url{indiamart.com} & Yagi antennas, bulk components \\
\textbf{SP Road, Bengaluru} & Physical & --- & Cables, connectors, SMA adapters \\
\bottomrule
\end{tabular}
\end{table}

% ============================================================================
\section{Recommended Shopping Carts}
% ============================================================================

\subsection{Budget Setup: RX-Only Ground Station (\rupee 7,800)}

\begin{table}[h!]
\centering
\begin{tabular}{@{}llll@{}}
\toprule
\textbf{Item} & \textbf{Price} & \textbf{Store} & \textbf{Purpose} \\
\midrule
RTL-SDR Blog V4 Kit & \rupee 6,000 & FabToLab / Robu.in & RX receiver \\
SX1276 LoRa Module (868~MHz) & \rupee 500 & Robu.in & LoRa TX uplink \\
868~MHz quarter-wave antenna & \rupee 300 & DIY / Robu.in & Signal capture \\
USB extension cable (3m) & \rupee 200 & Amazon.in & Reduce PC noise \\
SMA adapters & \rupee 300 & SP Road / Amazon & Connectivity \\
RG58 coax cable (3m) & \rupee 500 & SP Road & Antenna feed \\
\midrule
\textbf{TOTAL} & \textbf{\rupee 7,800} & & \\
\bottomrule
\end{tabular}
\end{table}

\subsection{Recommended Setup: HackRF TX/RX Station (\rupee 28,816)}

\begin{table}[h!]
\centering
\begin{tabular}{@{}llll@{}}
\toprule
\textbf{Item} & \textbf{Price} & \textbf{Store} & \textbf{Purpose} \\
\midrule
HackRF One V1.9.1 + Case + TCXO & \rupee 25,744 & FabToLab (Opt 85) & TX+RX 1--6~GHz \\
ANT700 antenna (300--1100~MHz) & \rupee 2,372 & FabToLab & Covers 868~MHz LoRa \\
SX1276 LoRa Module (868~MHz) & \rupee 500 & Robu.in & Dedicated LoRa TX \\
USB extension cable & \rupee 200 & Amazon.in & Reduce noise \\
\midrule
\textbf{TOTAL} & \textbf{\rupee 28,816} & & \\
\bottomrule
\end{tabular}
\end{table}

\subsection{Premium Setup: Portable Standalone Station (\rupee 60,759)}

\begin{table}[h!]
\centering
\begin{tabular}{@{}llll@{}}
\toprule
\textbf{Item} & \textbf{Price} & \textbf{Store} & \textbf{Purpose} \\
\midrule
HackRF One V1.9.1 + Case + TCXO & \rupee 25,744 & FabToLab (Opt 85) & Core SDR \\
PortaPack H4M & \rupee 29,503 & FabToLab & Standalone (no PC!) \\
ANT700 (300--1100~MHz) & \rupee 2,372 & FabToLab & LoRa band \\
800--2200~MHz antenna & \rupee 3,140 & FabToLab & Wideband satellite \\
\midrule
\textbf{TOTAL} & \textbf{\rupee 60,759} & & \\
\bottomrule
\end{tabular}
\end{table}

\subsection{Research Setup: Full-Duplex Satellite Lab (\rupee 1,05,000)}

\begin{table}[h!]
\centering
\begin{tabular}{@{}llll@{}}
\toprule
\textbf{Item} & \textbf{Price} & \textbf{Store} & \textbf{Purpose} \\
\midrule
LimeSDR Mini 2.0 & \rupee 40,000 & TheEngineerStore & Full-duplex TX+RX \\
ADALM-Pluto & \rupee 25,000 & element14 India & Secondary SDR / MIMO \\
868~MHz 12~dBi Yagi antenna & \rupee 2,000 & IndiaMART & Directional satellite \\
Wideband LNA (SPF5189Z) & \rupee 2,500 & FabToLab & Amplify weak signals \\
Raspberry Pi 4 (4GB) & \rupee 5,000 & Robu.in & 24/7 station computer \\
SX1262 LoRa module & \rupee 1,200 & ElectronicsComp & Advanced LoRa TX \\
Cables, adapters, enclosure & \rupee 3,000 & SP Road / Amazon & Infrastructure \\
\midrule
\textbf{TOTAL} & \textbf{$\sim$\rupee 78,700} & & \\
\bottomrule
\end{tabular}
\end{table}

% ============================================================================
\section{Import and Regulatory Notes}
% ============================================================================

\subsection{Customs and Compliance}

\begin{itemize}
    \item \textbf{WPC Approval:} Radio transmitting equipment in India may require 
          Wireless Planning and Coordination (WPC) approval. SDR receivers are 
          generally exempt; transmitters may need licensing depending on frequency and power.
    \item \textbf{BIS Certification:} Bureau of Indian Standards registration may be 
          required for direct international imports. Buying from Indian retailers avoids this.
    \item \textbf{ISM Band:} Transmission at 868~MHz below 25~mW does NOT require a 
          license in India under ISM band regulations.
    \item \textbf{GST:} Some Indian electronics sites show prices \textbf{excluding 18\% GST}. 
          Always verify the final checkout price.
    \item \textbf{Recommendation:} Buy from Indian stores (FabToLab, element14, Mouser.in) 
          to avoid customs detention and compliance headaches.
\end{itemize}

\subsection{Legal Use of SDR in India}

\begin{table}[h!]
\centering
\caption{Legal Status of SDR Activities in India}
\begin{tabular}{@{}lll@{}}
\toprule
\textbf{Activity} & \textbf{Legal?} & \textbf{Notes} \\
\midrule
Receiving any radio signal & Yes & No license needed to listen \\
Transmitting on ISM 868~MHz ($<$25~mW) & Yes & No license needed \\
Transmitting on amateur bands & License & Requires ASOC amateur license \\
Receiving weather satellites & Yes & Public broadcast, free \\
ADS-B reception & Yes & Public broadcast \\
GPS reception & Yes & Public broadcast \\
Intercepting cell phone signals & \textbf{No} & Illegal under Indian Telegraph Act \\
Jamming any signal & \textbf{No} & Serious criminal offense \\
High-power TX ($>$1~W) & License & Requires WPC approval \\
\bottomrule
\end{tabular}
\end{table}

% ============================================================================
\section{Conclusion}
% ============================================================================

This guide covers the complete spectrum of SDR hardware available in India as of 
June 2026, from the entry-level RTL-SDR V4 at \rupee 6,000 to the research-grade 
USRP B210 at \rupee 2,50,000.

\subsection{Summary of Options}

\begin{table}[h!]
\centering
\begin{tabular}{@{}llll@{}}
\toprule
\textbf{Budget} & \textbf{Best Choice} & \textbf{Price} & \textbf{Capability} \\
\midrule
Starter & RTL-SDR V4 & \rupee 6,000 & RX only, learn basics \\
Serious & HackRF One + TCXO & \rupee 25,744 & TX+RX, 1--6~GHz \\
Portable & HackRF + PortaPack & \rupee 60,759 & Standalone, no PC \\
Research & LimeSDR Mini 2.0 & \rupee 40,000 & Full-duplex, SatNOGS \\
Professional & BladeRF 2.0 xA9 & \rupee 75,000 & MIMO, 61~MHz BW \\
ISRO-grade & USRP B210 & \rupee 2,50,000 & Gold standard \\
\bottomrule
\end{tabular}
\end{table}

For the OpenOrbitLink project, the \textbf{HackRF One V1.9.1 with TCXO from FabToLab 
at \rupee 25,744} represents the optimal balance of cost, capability, and availability. 
It provides transmit capability for LoRa uplinks, covers the entire 1~MHz to 6~GHz 
spectrum, and is in stock for immediate delivery in India.

Combined with OpenOrbitLink's Android app and Python backend, this creates a 
\textbf{complete, open-source satellite communication system} for under \rupee 30,000 
--- compared to Starlink's \rupee 4,10,000 over 3 years.

\vspace{0.5cm}
\begin{center}
    \textit{``Making satellite connectivity accessible to everyone, everywhere.''}
    
    \textbf{OpenOrbitLink --- Affordable SDR, Zero Subscription, Infinite Possibilities.}
    
    \vspace{0.3cm}
    \small{GitHub: \url{https://github.com/Ashutosh0x/OpenOrbitLink}}
\end{center}

\end{document}
